% Terminal apparatus of the concise study: the scope note and the References
% for the sources this companion actually uses, at the loci it cites. The
% revision timestamp and the rights colophon follow from the shared
% generation record.

\runninghead{Scope and References}
\section{Appendix: Scope and Qualifications}

{\footnotesize
\textbf{Edition, Mass and occurrence.} The Mass treated is \latin{Dominica
prima Adventus}, I classis, \latin{Statio ad S. Mariam maiorem}, of the
\work{Missale Romanum}, editio typica, Typis Polyglottis Vaticanis 1962,
printed pp.~1--2, nos.~1--10, as said on Sunday 29~November 2026 in the
universal 1962 calendar. The Sunday, its first class, the omission of the
commemoration of St Saturninus, the colour, the omission of the Gloria, the
Credo and the Trinity Preface rest on the Missal's general rubrics
(\work{Rubricae generales}, abbreviated RG in the note under the map)
and its printed directions, as \texttt{research/context.md} records. The Ordinary and
the Office are not treated.

\textbf{Text and English.} The Latin of the ten texts was read on page images
of the CMAA facsimile of the 1962 typical edition and collated with the
Ratisbon printing of 1862 (Pustet); the collation is in
\texttt{propers/verified.md}, and the expansive study prints every text whole.
The lessons and chants are quoted in the Douay--Rheims in Challoner's
revision; the three orations and the Communion in the Lasance \work{New Roman
Missal} (Benziger, 1945), public domain in the United States by non-renewal.
The Introit and Offertory antiphon is given in the Douay without its ``O
Lord,'' which the Missal's \latin{Ad te levávi} lacks. None of this English is an approved liturgical translation of the 1962
Missal. Where the argument turns on the Latin of a chant or an oration, single
words or a clause are rendered literally in the sentence; these renderings are
not a version of the prayer. Chrysostom, Cyril of Jerusalem, Augustine on
Psalm~24 and the \work{Confessions} are quoted in the English of the Nicene
and Post-Nicene Fathers; Cyril of Alexandria in R.~Payne Smith's English
(1859); Bellarmine in John O'Sullivan's abridged English (1866); Schuster in
the English of 1924. Wherever a Father's own Latin is quoted, the English
beside it is the rendering made for the expansive study.

\textbf{Reception.} The four interpretations are developed in full in the
expansive study; their loci, editions and bounds are recorded in
\texttt{research/scope.md} and \texttt{research/interpretations.md}. Every
connection drawn between two texts of the Mass belongs to the interpretations
unless an author is named as making it; no Father cited here comments on this
Mass as a whole. Gregory, Ambrose, Cyril of Alexandria, Augustine's
Letter~199 and Bede comment on Luke's discourse; Chrysostom and Cyril of
Jerusalem on the parallel in Matthew. Gregory, Ambrose, Bede and Augustine's
\work{Enarrationes} on Psalm~84 were read in Migne's Latin as transcribed on
Wikisource and in the Corpus Corporum, not collated with page images;
Augustine's Letter~199 in the augustinus.it text; the Latin of the
\work{Confessions} in J.~J. O'Donnell's text as given by The Latin Library;
Aquinas on Romans and on Psalm~24 in the Corpus Thomisticum text; Rupert,
Sicard and Durandus in uncorrected optical text, whose quotations are
normalised in spelling.
Rupert of Deutz, Sicard of Cremona and William Durandus are ecclesiastical
writers; they and Schuster, who is among the Blessed, are cited in support,
for what they say of texts this Mass reads. Cassiodorus on the psalms, Jerome
on Matthew, Augustine's \work{City of God}~XX and the \work{Catena aurea} were
not read.

\textbf{Chronology.} Every biblical date on page~2 is generated from the
repository's Scripture chronology corpus; the expansive study's Scriptural
Date and Location appendix gives each passage's sources and evidence.

\textbf{Rights and use.} The Latin is published on 17 U.S.C. 103(b) together
with its public-domain witness, the Pustet printing of 1862. The Douay--Rheims
and the Lasance English are quoted from registered editions. The other
quotations are short and come from editions published before 1931, named
below, or from ancient and medieval Latin. Scripture, liturgical texts and
quoted translations are not relicensed by the repository's licence. This is
a study aid to the 1962 Missal, not an official liturgical book, translation
or critical edition.
\par}

\section{References}

{\footnotesize
\begin{itemize}[leftmargin=1.2em,itemsep=.15em]
\item \work{Missale Romanum ex decreto Sacrosancti Concilii Tridentini
restitutum, Summorum Pontificum cura recognitum}, editio typica (Typis
Polyglottis Vaticanis, 1962), CMAA facsimile: \work{Rubricae generales};
Proprium de Tempore, pp.~1--2, nos.~1--10.
\item \work{Missale Romanum} (Ratisbon: Pustet, 1862), pp.~1--3.
\item Lasance, F.~X., and Walsh, F.~A., \work{The New Roman Missal in Latin and
English}, revised edition (New York: Benziger Brothers, 1945), ``First Sunday
of Advent,'' pp.~103--106.
\item \work{The Holy Bible Translated from the Latin Vulgate} (Douay--Rheims,
Challoner revision), Project Gutenberg text; \work{Biblia Sacra Vulgatae
Editionis} (Clementine Vulgate), for the comparisons with the Missal's forms.
\item Ambrose, \work{Expositio evangelii secundum Lucam} X.37--42, Migne
Latin text, PL~15 (Wikisource).
\item Augustine, \work{Enarrationes in Psalmos} 24.1--3 and 24.10, NPNF,
first series, vol.~8; 84.9 and 84.15, Migne Latin text (Wikisource);
\work{Epistula} 199 to Hesychius, augustinus.it Latin text;
\work{Confessions} VIII.5.12 and VIII.12.29, trans.\ J.~G. Pilkington, NPNF,
first series, vol.~1, and VIII.5.12, Latin text ed.\ J.~J. O'Donnell (The
Latin Library).
\item Bede, \work{In Lucae evangelium expositio} VI, on 21:32--33, Migne Latin
text (Corpus Corporum).
\item Bellarmine, Robert, \work{Explanatio in Psalmos}, on Ps~24:2--4 and
Ps~84:8 and 13, trans.\ John O'Sullivan, \work{A Commentary on the Book of
Psalms} (Dublin: Duffy, 1866).
\item Chrysostom, John, \work{Homilies on Romans} 24, NPNF, first series,
vol.~11; \work{Homilies on Matthew} 77, NPNF, first series, vol.~10.
\item Cyril of Alexandria, \work{Commentary on the Gospel of Saint Luke},
sermon 139, trans.\ R.~Payne Smith (Oxford, 1859).
\item Cyril of Jerusalem, \work{Catechetical Lectures} 15.1, 15.3--4, NPNF,
second series, vol.~7.
\item Gregory the Great, \work{Homiliae in Evangelia} I.1, nn.~1--6, Migne
Latin text (Wikisource).
\item Thomas Aquinas, \work{Super Epistolam ad Romanos}, cap.~13, lect.~3,
and \work{Super Psalmos}, on Ps~24, nn.~1 and 4, Corpus Thomisticum text.
\item Durandus, William, \work{Rationale divinorum officiorum} VI, on the
first Sunday of Advent (Lyon, 1612).
\item Rupert of Deutz, \work{De divinis officiis} III.2, PL~170.
\item Schuster, Ildefonso, \work{The Sacramentary (Liber Sacramentorum)},
vol.~1 (London: Burns, Oates and Washbourne, 1924), ``First Sunday of
Advent,'' pp.~321--323.
\item Sicard of Cremona, \work{Mitrale} V, on the first Sunday of Advent,
PL~213.
\end{itemize}
\par}
